\section{Further research questions}
\label{future:sec}

The formulas above suggest concrete next problems. The questions below
are proposals, not assertions of currently unproved identities.
They are ordered roughly by how directly the present proofs provide
the next calculation.

\subsection{Cubic and higher coincident products}

\paragraph{1. Reduce the third diagonal Tornheim derivative.}
Theorem~\ref{cub:diagonal} isolates $\omega'''(0)$ as the remaining
global coordinate for the cubic digamma finite part. Seek an exact
reduction of this derivative to a stated algebra of ordinary zeta
derivatives, Gamma or Barnes values, and convergent logarithmic
integrals. A useful first target is a second independent functional
identity for the cubic integral, rather than an unqualified
integer-relation search in a large constant basket.
If no such reduction is found, an explicit minimal integral
representation is still valuable, but it must not be called an
independence result.

\paragraph{2. Determine every asymmetric cubic ray derivative.}
In the local decomposition \eqref{sp:meromorphicH}, the next Taylor
layer involves $H_{AA}=H_{BB}$, $H_{AB}$, $H_{AC}=H_{BC}$, and
$H_{CC}$. The exact $C$ axis determines $H_{CC}$.
The diagonal cubic coefficient supplies one linear combination of
the other three. Can two additional exact restrictions or functional
equations determine the remainder? A concrete deliverable would be
a formula for $W_{a,b,c}'''(0)$ with its smallest explicit collection
of convergent-integral coordinates. This would clarify exactly where
ordinary single-zeta data cease to determine all directions.

\paragraph{3. Give the explicit global kernel at degree four.}
The local subtraction at any number of factors is already finite:
\eqref{cub:subset-numerator} lists the needed Taylor coefficients.
The unsolved part is the arithmetic of the global Fourier
zero-frequency constraint. At degree four, separate the sign
partitions of $n_1+n_2+n_3+n_4=0$ and identify a finite set of
constrained multiple sums whose meromorphic continuations represent
all moments. The first useful test is the fully coincident
$\FP\int_0^1\psi(x)^4\,dx$ with every counterterm written out.
This would extend the present cubic kernel without assuming
single-zeta closure.

\paragraph{4. Analyze triple collisions on unequal grids.}
The incoming pair formulas allow unequal spatial dilations.
A triple product introduces several collision loci and a possible
full three-way coincidence. Derive a chamber formula in the shift
variables, identify the exact changes on crossing each collision
wall, and compare the successive pair-collision limits with the
direct triple-coincident finite part. The challenge is to preserve
one stated coordinate prescription throughout; repeated finite-part
operations need not commute automatically.

\subsection{Contact distributions and coordinate transport}

\paragraph{5. Extend the contact generator to several shifts.}
Theorem~\ref{cont:main} determines the entire point-supported
correction for two factors and one shift. For a three-factor
correlation, the collision sets are lines and points in a two-shift
torus. Determine the distributions supported on those sets and their
intersections. The pairwise terms should recover
\eqref{cont:E}; the new information is the correction at a joint
intersection. A meaningful answer should specify test-function
normalizations and not introduce a product of singular
distributions without a definition.

\paragraph{6. Characterize invariant products exactly.}
Corollary~\ref{invg:stieltjes-products} gives a clean sufficient
criterion: at most one zero Stieltjes index. It is not a complete
classification of products or linear combinations of products.
Apply Theorem~\ref{invg:classification} to the finite Taylor jets of
the analytic factors and identify all cancellations among the
forbidden monomials. One concrete starting point is the span of
$\gamma_i\gamma_j\gamma_k$ at fixed total index, with coefficients
in a prescribed constant field. The field matters because a
functional cancellation and an arithmetic relation are different
questions.

\paragraph{7. Determine the coordinate-defect composition law.}
For two tangent coordinates $f$ and $g$, express the defect for
$f\circ g$ in terms of the defects for $f$ and $g$ and the
corresponding pullbacks. Identify a finite representation of this
law on pole order at most $R$ and logarithmic degree at most $L$.
The invariant subspace and its exact codimension are now known.
The remaining structural question is the action on the quotient,
including whether a natural coordinate choice gives a canonical
representative for each obstruction class.

\subsection{Harmonic Laurent data and primitives}

\paragraph{8. Move all the harmonic exponents independently.}
Theorem~\ref{h2:thm:direction} treats an outer exponent and the first
inner exponent while every trailing exponent remains one.
At the first new case
$\zeta_a(s,t_1,t_2)$, perturb all three orders independently.
Find the complete set of local polar divisors, then determine which
principal coefficients and finite parts admit finite Gamma,
Bernoulli, and Stirling formulas without an unevaluated
holomorphic remainder. General multiple-zeta Laurent algorithms
already exist \cite{h2:MOW}; the target is a compact explicit
specialization that preserves the common shift and all directions.

\paragraph{9. Classify directions with explicit higher regular terms.}
The two-exponent theorem stops at the finite part for a reason:
derivatives of its holomorphic remainder generally enter later.
The exact normal slice in Theorem~\ref{h2:thm:slice} is an exception
with an all-order Stieltjes formula. Determine other integer
outer-order slices or direction combinations for which the
remainder cancels to a prescribed order. A finite criterion would
organize which higher coefficients genuinely require new multiple
zeta derivatives.

\paragraph{10. Normalize primitives using Gamma multiplication.}
Equation~\eqref{h2:eq:primitive} supplies a finite primitive in the
shift variable, and \eqref{cont:Bernoulli-inverse} supplies
mean-zero periodic primitives. Study their compatibility with
reflection, multiplication of the shift, and rational shifts.
At low derivative orders, compare the resulting constants with
standard Barnes multiple-Gamma normalizations. Any claimed
simplification should specify both the additive constant and the
branch of the logarithm.

\subsection{Colored sums and exact Gaussian certificates}

\paragraph{11. Reduce mixed inner signs at arbitrary depth.}
The rational theorem settles the sector with all inner orders
nonpositive and arbitrary positive integer weights. Positive
integer orders at depth two are handled by partial fractions.
Develop a finite nested-polylogarithm reduction for arbitrary depth
with both positive and nonpositive inner orders, while retaining
the full complex outer order. Coincident colors and repeated poles
must be built into the formula. An affine shift alone is a routine
Lerch or Hurwitz extension and is less informative than this mixed
sector.

\paragraph{12. Detect arithmetic cancellation of depth-two jets.}
At roots of unity, \eqref{sp:DFTjets} and
\eqref{sp:DFTstieltjes} completely describe the depth-one terms.
Which character projections, symmetrizations, or distribution
combinations cancel the residual depth-two order derivatives in
\eqref{sp:11}? A useful result would be an exact finite linear map
whose kernel identifies combinations reducible to Gamma,
Hurwitz-zeta, and generalized-Stieltjes data. Such a map should
retain the allowed derivative directions instead of specializing
before differentiation.

\paragraph{13. Seek exact $S_6$ and $S_8$ certificates.}
Use the frozen target vectors at the pinned revision and enlarge the
proved relation space, not the numerical precision alone. The
existing quotient obstruction identifies relations that cannot
suffice; the verified $S_4$ octahedral proof identifies additional
families worth testing. For every added regularized relation,
prove the regulator prescription and retain its exact rational
coefficient row. A successful certificate must sum to the target
identity exactly. If a restricted search fails, publish the
precisely defined relation space and its separating functional,
without extending the impossibility statement beyond that space.

\paragraph{14. Formalize the finite proof kernels.}
The coefficient extractions in the subset completion, contact
generator, and harmonic residue polynomial are finite algebra.
Separate that algebra from the analytic lemmas of normal
convergence and meromorphic continuation, and formalize each part
with explicit interfaces. In particular, a reusable proof of the
Taylor-subtraction lemma in several spectral variables would
prevent accidental interchange of a coordinate finite part and a
spectral Laurent constant in future integrations.
